\chapter{Nói Với Lớp Trưởng, Cán Sự}
\chapterintro{Không để người giữ nề nếp thành người gánh hết}{Lớp trưởng và cán sự rất cần được nói đúng cách để đỡ áp lực mà vẫn có vai trò.}

\section{Em không cần làm thay cả lớp}
\begin{manthoai}{Màn thoại}
\textbf{Thầy/Cô:} Em không cần làm thay cả lớp, em chỉ cần giúp lớp đúng việc.\\
\textbf{Lớp trưởng:} Dạ.
\end{manthoai}
\begin{dungkhinào}
Hợp khi lớp trưởng đang ôm quá nhiều việc hoặc mang tâm lý phải gánh mọi thứ.
\end{dungkhinào}
\begin{nhipthem}
Câu này giúp hạ áp lực và cũng chỉnh lại vai trò cho đúng.
\end{nhipthem}
\ngancach

\section{Em nhắc lớp bằng giọng bạn học trước, chưa cần giọng quản lý}
\begin{manthoai}{Màn thoại}
\textbf{Thầy/Cô:} Em nhắc lớp bằng giọng bạn học trước, chưa cần giọng quản lý.\\
\textbf{Cán sự:} Dạ em hiểu.
\end{manthoai}
\begin{dungkhinào}
Dùng khi cán sự bắt đầu nhắc bạn theo kiểu quá gắt, làm phản tác dụng.
\end{dungkhinào}
\begin{nhipthem}
Đây là cách góp ý rất hợp cho giờ sinh hoạt hoặc trao đổi riêng sau lớp.
\end{nhipthem}
\ngancach

\section{Giữ lớp không phải là quát to hơn lớp}
\begin{manthoai}{Màn thoại}
\textbf{Thầy/Cô:} Giữ lớp không phải là quát to hơn lớp.\\
\textbf{Lớp trưởng:} Dạ.
\end{manthoai}
\begin{dungkhinào}
Hợp khi muốn dạy học sinh nòng cốt về cách giữ tập thể mà không biến thành cuộc đấu âm lượng.
\end{dungkhinào}
\begin{nhipthem}
Có thể nối bằng việc tập cho lớp trưởng một câu gọi lớp thật ngắn và rõ.
\end{nhipthem}
\ngancach

\section{Em giữ nhịp thôi, đừng giữ hết việc}
\begin{manthoai}{Màn thoại}
\textbf{Thầy/Cô:} Em giữ nhịp thôi, đừng giữ hết việc.\\
\textbf{Lớp trưởng:} Dạ.
\end{manthoai}
\begin{dungkhinào}
Hợp khi lớp trưởng quá có trách nhiệm tới mức ôm cả phần của các bạn khác.
\end{dungkhinào}
\begin{nhipthem}
Câu này giúp em ấy hiểu vai trò là điều phối, không phải gánh thay tập thể.
\end{nhipthem}
\ngancach

\section{Nhắc lớp xong nhớ cảm ơn lớp một câu}
\begin{manthoai}{Màn thoại}
\textbf{Thầy/Cô:} Nhắc lớp xong nhớ cảm ơn lớp một câu, để bạn nghe em như nghe người cùng đội.\\
\textbf{Cán sự:} Dạ.
\end{manthoai}
\begin{dungkhinào}
Dùng khi cán sự làm đúng việc nhưng giọng còn cứng, dễ gây phản ứng ngược.
\end{dungkhinào}
\begin{nhipthem}
Chỉ cần thêm một câu cảm ơn nhỏ, cách nhắc của học sinh nòng cốt đã mềm đi rất nhiều.
\end{nhipthem}
